\section{Linear Models}
Using inherently interpretable algorithms.
\subsection{Linear Regression}
Predicts target as weighted sum of inputs.
\begin{equation}
	y = \beta_0 + \beta_1 x_1 + \dots + \beta_p x_p + \epsilon
	\label{eq:linear_regression}
\end{equation}
Estimated usually by Ordinary Least Squares (OLS)
\begin{equation}
	\hat \beta = \argmin{\beta_0, \dots , \beta_p} \sum_{i=1}^n \lr{y^{(i)} - \lr{\beta_0 + \sum_{j=1}^p}\beta_jx_j^{(i)}}²
	\label{eq:OLS}
\end{equation}
Each weight represents the effect of one feature, holding others constant.
We can compute \textbf{confidence intervals} for each coefficient.
\paragraph{Regression Output}
\begin{table}[H]
	\caption{Regression Output}\label{tab:regression_output}
	\begin{center}
		\begin{tabular}{l|r|r|r}
			\hline
			Feature     & Weight$\hat \beta_j$ & SE  & $|t|$ \\
			\hline
			Temperature & 110.7                & 7.0 & 15.7  \\
			Humidity    & -17.4                & 3.2 & 5.5   \\
			\hline
		\end{tabular}
	\end{center}
\end{table}
\begin{itemize}
	\item Weight: Estimated influence of $x_i$ over the outcome $y$
	\item SE (Standard Error): Uncertainty (variance) of the estimated coefficient
	\item $|t|$ t-statistic: Strength of evidence that $\beta_j \neq 0$
\end{itemize}
\begin{equation}
	t_{\beta_j} = \frac{\hat \beta_j}{SE(\hat \beta_j)}
	\label{eq:t-statistic}
\end{equation}
Large, certain weights = important features. Small or uncertain weights = weak evidence. Features wit highest $|t|$ values are most significant.

\begin{itemize}
	\item Effect Plot: $effect_j^i = \beta_j \cdot x_j^i$ Contribution of feature j to prediction for instance i.
	\item Different encodings for variables.
\end{itemize}


\paragraph{Notes on Linear Regression}
\begin{itemize}
	\item Perfectly understood
	\item Only captures linear relationships
	\item Ignores feature interactions, must be added manually ($x_1x_2, x²$ etc.)
	\item Lower accuracy compared to more flexible Models.
\end{itemize}


\subsection{Lasso - Least Absolute Shrinkage and Selection Operator}
Using models that keep only the relevant features helps with interpretability.
\begin{itemize}
	\item Adds L1 Norm, Performs feature selection, regularization
	\item Enforces sparsity in Linear model
\end{itemize}
\begin{equation}
	\underset{\beta}{min} \frac{1}{n}\sum:{i=1}^n\lr{y^{(i)}-x_i^T \beta }² + \lambda \norm{\beta}_1
	\label{eq:LASSO-L1}
\end{equation}
As $\lambda$ gets increased, only the most important features remain. This increases interpretability.



\subsection{Logistic Regression}
\begin{itemize}
	\item Classification problems with two classes.
	\item Models Probability of one class given the features.
	\item Extension of linear regression for classification.
	\item Uses the output of a linear regression and maps it to a Probability.
\end{itemize}
\paragraph{Logistic / Sigmoid function}
\begin{equation}
	logistic\lr{\eta} = \frac{1}{1+ \exp{-\eta}}
	\label{eq:Log-sigmoid-fun}
\end{equation}

\paragraph{Logistic Model}
\begin{equation}
	P\lr{y^{(i)} = 1}  = \frac{1}{1 + \exp (- (\beta_0 + \beta_1 x_1^{(i)} + \dots + \beta_p x_p^{(i)}))}
	\label{eq:logisticregression}
\end{equation}
\begin{itemize}
	\item Linear predictor is transformed via the logistic function.
	\item Coefficients do not act linearly on the probability. Interpretation is happens on the log-odds scale.
	\item A one unit increase in $x_j$ multiplies the odds of $y=1$ vs. $y=0$ by a factor of $\exp(\beta_j)$
	\item Predicts probabilities, which allows for better decision making. Pred with 99\% probability is more confident than one with 51\%.
	\item Relationships between features and outcomes remain transparent.
\end{itemize}

\subsection{Linear Regression Extensions}
Linear Regression Models can be extended to handle some things that normal linear Regression can not.
\begin{itemize}
	\item $y$ distribution is not Gaussian: Use Generalized Linear Models (GLMs)
	\item The features interact: Add interactions manually
	\item The relationship between $x$ and $y$ is not linear: Use Generalized Additive Models (GAMs)
\end{itemize}

\subsubsection{Generalized Linear Models}
Use this if the output $y$ follows a Non-Gaussian distribution.
For example if we want to predict counts. They follow a Poisson distribution (discrete values, $Var(y) = \E{y}$)
\begin{equation}
	g(\E{Y|x}) = \beta_0 + \beta_1 +  \dots + \beta_p x_p = x^T\beta
	\label{eq:GLM-formula}
\end{equation}
It has three components

\begin{itemize}
	\item \textbf{Random Component}: A distribution of the exponential family that defines $\mathbb{E}_Y$
	\item \textbf{Systematic component}: Linear predictor $\eta = X^T\beta$. (LR Term)
	\item \textbf{Link function}: Function $g(\cdot)$ that links the mean to $\eta$
\end{itemize}

Our models prediction is: $\mathbb{E}[y] = g^{-1}(x^T\beta)$ which could be 7.6 for example customer count.
But with GLM, since we have defined the distribution of y (Poisson), we can make several inferences using the predicted mean.
Of course this means, that we have to choose $\E{y}$ correctly, according to our data, or the inferences will be wrong.
\begin{itemize}
	\item Probability that we get more than 15 customers, when 7.6 was predicted as mean.
	\item Catch outliers
	\item Calculate the variance of our prediction.
\end{itemize}
\paragraph{Distributions and their canonical link functions}
\begin{itemize}
	\item Poisson: Log: $g(\mu) = \eta$
	\item Binomial: Logit: $\ln\lr{\frac{\mu}{1-\mu}}=\eta$
	\item Normal: Identity: $\mu=\eta$
\end{itemize}

\textit{GLMs transform model outputs $y$}

\subsubsection{Feature Interactions}
\begin{itemize}
	\item Features interact with each other.
	\item E.g. Nr of bike Rentals: Temperature and Working/Non-Working day. on days of people only rent bikes when its warm. On working days, people ride to work regardless of weather.
	\item In LR we need to model these interactions explicitly. Add Interaction terms (Product of two features)
	\item Modelling of interactions between different feature types:
	      \begin{itemize}
		      \item \texttt{Cat} x \texttt{Num}: Existence of category x Num feature value.
		      \item \texttt{Cat} x \texttt{Cat}: \texttt{AND} between one hot encoded features.
		      \item \texttt{Num} x \texttt{Num}: Multiply the features
	      \end{itemize}
\end{itemize}


\subsubsection{Generalized Additive Models}
In a linear model, the increase of a feature by 1 has the same effect everywhere (over the whole range of that features values)
But there is not always a linear relationship over the entire range. (Growth, saturation, decline). We need to model these nonlinear relationships.

E.g. 10 to 11 degrees, bike rentals increase, but 40 to 41 degrees they decline (to hot).
\paragraph{Approaches to Non-Linearity}
\begin{itemize}
	\item Replace features with transformed versions. $log(x), \sqrt{x}, x²$
	\item Discretizise categorical features into Bins/Quantiles (one hot encoded) $\to$ model becomes step function.
	\item GAMs: Let model learn smooth non-linear functions of features.
\end{itemize}

\paragraph{GAMs}
\begin{equation}
	g(\Ewrp{Y}{y|x}) = \beta_-1 + f_1(x_1) + f_2(x_2) + \dots + f_p(x_p)
	\label{eq:GAM}
\end{equation}
\begin{itemize}
	\item GAMs are GLMs but the features are transformed non-linearly with a function.
	\item GAMs allow GLMs to learn nonlinear relationships.
	\item $f_i(x_i)$ is usually a learned spline function.
	      \begin{itemize}
		      \item Think of them as $k$ new features that are obtained via $k$ spline functions. The spline functions outputs combined via a weighted sum model the non-linearity.
		      \item The range is split into $k$ sections/knots (e.g. by percentile)
		      \item For each section, a spline is fit e.g cubic $y = ax³ + bx² + cx + d$
		      \item Then the outputs of these splines is combined via a weighted sum.
	      \end{itemize}
\end{itemize}

Maintain interpretability, because the features are still additive across features.


\textit{GAMs transform model inputs $x$ with learned non-linear functions.}


